\begin{frame}[fragile]{Backup: Python Implementation of RAII}
	\begin{figure}
		\centering
		\begin{columns}[T,onlytextwidth]
			\begin{column}{0.5\textwidth}
				\begin{minted}{py}
def read_file():
    f = open("data.txt", "r")
    content = f.read()
    return content
		\end{minted}
				\rule{\linewidth}{0.5mm}
				\begin{minted}{py}
def read_file():
    with open("data.txt", "r") as f:
        content = f.read()
    return content
				\end{minted}
			\end{column}
			\begin{column}{0.5\textwidth}
				\begin{minted}{py}
class FileReader:
    def __init__(self, path):
        self.file = open(path, "r")
    def read(self):
        return self.file.read()
    def __del__(self):
        if not self.file.closed:
            self.file.close()
def read_file():
    reader = FileReader("data.txt")
    return reader.read()
				\end{minted}
			\end{column}
		\end{columns}
		\caption{Example: RAII concepts in Python}
	\end{figure}
\end{frame}

\begin{frame}[fragile]{Backup: C++ Lock Guard}
	\begin{figure}
		\centering
		\begin{columns}[T,onlytextwidth]
			\begin{column}{0.5\textwidth}
				\begin{minted}{cpp}
#include <mutex>
std::mutex m;
void critical_section() {
	m.lock();
	// critical work
	m.unlock();
} // unlock may not be executed
  // on exception throw
				\end{minted}
			\end{column}
			\begin{column}{0.5\textwidth}
				\begin{minted}{cpp}
#include <mutex>
std::mutex m;
void critical_section() {
    std::lock_guard<std::mutex> lock(m);
    // critical work
} // mutex is automatically unlocked
				\end{minted}
			\end{column}
		\end{columns}
		\caption{Example: C++ Lock Guard for concurrency control}
	\end{figure}
\end{frame}

\begin{frame}[fragile]{Backup: C++ Shared Pointer (Reference Counting)}
	\begin{figure}
		\centering
		\begin{minted}{cpp}
#include <iostream>
#include <memory>
int main() {
    std::shared_ptr<int> a = std::make_shared<int>(42);
    {
		// share/copy pointer
        std::shared_ptr<int> b = a; // reference count increases
        std::cout << *b << '\n';
    } // b goes out of scope, reference count decreases
    std::cout << *a << '\n';
} // a goes out of scope, resource is released
		\end{minted}
		\caption{Example: C++ Reference Counting}
	\end{figure}
\end{frame}

\begin{frame}[fragile]{Backup: Use Valgrind to detect Memory Management problems (i)}
	\begin{figure}
		\centering
		\begin{minted}{cpp}
void leak() {
    int* data = new int[10];  // allocated
    data[0] = 42;
    // missing delete[] → leak
}
int main() {
    leak();
}
		\end{minted}
		\begin{minted}{sh}
g++ -g -O0 leak.cpp -o leak
valgrind --leak-check=full --show-leak-kinds=all ./leak
		\end{minted}
		\caption{Example: Use valgrind to find Memory Leaks in e.g. C++ (i)}
	\end{figure}
\end{frame}

\begin{frame}[fragile]{Backup: Use Valgrind to detect Memory Management problems (ii)}
	\begin{figure}
		\centering
		\begin{minted}[fontsize=\scriptsize]{text}
==122274== Memcheck, a memory error detector
==122274== Copyright (C) 2002-2024, and GNU GPL'd, by Julian Seward et al.
==122274== Using Valgrind-3.25.1 and LibVEX; rerun with -h for copyright info
==122274== Command: ./leak
==122274== HEAP SUMMARY:
==122274==     in use at exit: 40 bytes in 1 blocks
==122274==   total heap usage: 2 allocs, 1 frees, 73,768 bytes allocated
==122274== 40 bytes in 1 blocks are definitely lost in loss record 1 of 1
==122274==    at 0x4850613: operator new[](unsigned long) (vg_replace_malloc.c:729)
==122274==    by 0x400114A: leak() (leak.cpp:3)
==122274==    by 0x4001164: main (leak.cpp:8)
==122274== 
==122274== LEAK SUMMARY:
==122274==    definitely lost: 40 bytes in 1 blocks
==122274==    indirectly lost: 0 bytes in 0 blocks
==122274==      possibly lost: 0 bytes in 0 blocks
==122274==    still reachable: 0 bytes in 0 blocks
==122274==         suppressed: 0 bytes in 0 blocks
==122274== 
==122274== For lists of detected and suppressed errors, rerun with: -s
==122274== ERROR SUMMARY: 1 errors from 1 contexts (suppressed: 0 from 0)
		\end{minted}
		\caption{Example: Use valgrind to find Memory Leaks in e.g. C++ (ii)}
	\end{figure}
\end{frame}
